\documentclass{article}
\usepackage[fontsize=14pt]{fontsize}
\usepackage{fontspec}
\setsansfont{CMU Sans Serif}%{Arial}
\setmainfont{CMU Serif}%{Times New Roman}
\setmonofont{CMU Typewriter Text}%{Consolas}
\defaultfontfeatures{Ligatures={TeX}}
\usepackage[math-style=TeX]{unicode-math}
\usepackage[english, russian, ukrainian]{babel}

\usepackage{physics}
\usepackage{cleveref}

\usepackage[%
	a4paper,%
	footskip=1cm,%
	headsep=0.3cm,%
	top=2cm, %поле сверху
	bottom=2cm, %поле снизу
	left=2cm, %поле ліворуч
	right=2cm, %поле праворуч
    ]{geometry}


\usepackage{amsmath}
\usepackage{mathtools}

\title{Виведення впіввідношення Беннета}
\date{}
\begin{document}
\maketitle

\begin{equation}\label{eq:original}
    r^2 \frac{dp}{dr} = -\frac{1}{\pi c^2} \frac{d}{dr} \left( \frac{1}{2} I_z^2 \right)
\end{equation}


Множимо обидві частини на $dr$ і переписуємо:
\begin{equation}
    r^2 dp = -\frac{1}{\pi c^2} d\left( \frac{1}{2} I_z^2 \right)
\end{equation}

Інтегруємо від $r = 0$ до $r = a$:
\begin{equation}
    \int_{0}^{a} r^2 dp = -\frac{1}{\pi c^2} \int_{0}^{a} d\left( \frac{1}{2} I_z^2 \right)
\end{equation}

Ліва частина інтегрується по частинах:
\begin{align}
    \int_{0}^{a} r^2 dp &= r^2 p \Big|_{0}^{a} - \int_{0}^{a} 2r p \, dr \nonumber \\
    &= - \int_{0}^{a} 2r p \, dr \quad \text{(маємо на увазі $r^2 p \big|_{0}^{a}$ = 0)}
\end{align}

Права частина обчислюється безпосередньо:
\begin{align}
    -\frac{1}{\pi c^2} \left( \frac{1}{2} I_z^2 \Big|_{0}^{a} \right)
    &= -\frac{1}{\pi c^2} \left( 0 - \frac{1}{2} I_0^2 \right) \nonumber \\
    &= \frac{I_0^2}{2\pi c^2} \quad \text{(при } I_z(0) = I_0,\ I_z(a) = 0\text{)}
\end{align}

Прирівнюємо результати:
\begin{equation}
    - \int_{0}^{a} 2r p \, dr = \frac{I_0^2}{2\pi c^2}
\end{equation}

Множимо обидві частини на $-2\pi c^2$ і отримуємо підсумковий вираз:
\begin{equation}
    I_0^2 = 2c^2 \int_{0}^{a} 2\pi r p(r) \, dr
\end{equation}

Вважаючи, що електронна \( T_e \) та йонна \( T_i \) температури в плазмовому циліндрі постійні,
отримуємо

\begin{equation}\label{eq:Bennett}
    I_0^2 = 2c^2 \left( T_e + T_i \right) N,
\end{equation}
де \( N \) --- число частинок на одиницю довжини (так зване погонне число частинок)
плазмового циліндра:

\begin{equation}
    N = \int_{0}^{a} 2\pi r n(r) \, dr.
\end{equation}

Співвідношення \eqref{eq:Bennett} називається \emph{співвідношенням Беннета}.

\begin{small}
При отриманні співвідношення Беннета ми припускали, що
температура плазми постійна. У реальній плазмі протікання струму
супроводжується виділенням джоулевого тепла і нагріванням плазми. Якщо цей
процес відбувається повільно, то умову рівноваги можна вважати
наближено справедливою. Тому якщо за постійного погонного числа частинок
температура зростатиме, то для забезпечення рівноваги необхідно
постійно збільшувати струм.
\end{small}



Радіальний розподіл тиску \( p(r) \) можна виразити через \( B_0(r) \). Перепишемо (5.18) у вигляді
\begin{equation}\label{eq:5.24}
    \dv{p}{r} = -\frac{c}{8\pi r^2} \dv{r} \left[ r^2 B^2(r) \right]
\end{equation}
і проінтегруємо від \( r = 0 \) до довільного радіуса \( r \):
\begin{equation}\label{eq:5.25}
    p(r) = p(0) - \frac{1}{8\pi} \int_{0}^{r} \frac{1}{r^2} \dv{r} \left[ r^2 B^2(r) \right] \dd{r}.
\end{equation}

Як і вище, покладемо \( p(a) = 0 \), тоді
\begin{equation}\label{eq:5.26}
    p(0) = \frac{1}{8\pi} \int_{0}^{a} \frac{1}{r^2} \dv{r} \left[ r^2 B^2(r) \right] \dd{r}.
\end{equation}

Підставляючи \eqref{eq:5.26} в \eqref{eq:5.25}, отримуємо шукану залежність:
\begin{equation}\label{eq:5.27}
    p(r) = \frac{1}{8\pi} \int_{r}^{a} \frac{1}{r^2} \dv{r} \left[ r^2 B^2(r) \right] \dd{r}.
\end{equation}

Середній тиск \( \bar{p} \) у $z$-пінчі визначається рівнянням
\begin{equation}\label{eq:5.28}
    \bar{p} = \frac{1}{\pi a^2} \int_{0}^{a} 2\pi r p(r) \dd{r} = -\frac{1}{a^2} \int_{0}^{a} r^2
    \dv{p}{r} \dd{r}.
\end{equation}




Підставляючи в \eqref{eq:5.28} \( \dv{p}{r} \) з \eqref{eq:5.24}, отримуємо шукану залежність:
\begin{equation}\label{eq:5.29}
    \bar{p} = \frac{B^2(a)}{8\pi} = \frac{4I_0^2}{8\pi c^2 a^2} = \frac{I_0^2}{2\pi c^2 a^2}.
\end{equation}

Рівняння \eqref{eq:Bennett} і \eqref{eq:5.29} є інтегральними умовами рівноваги z-пінча. Вони не
залежать
від радіального розподілу струму. В експериментальних умовах розподіл струму визначається способом
утворення плазми та її провідністю.

В імпульсних розрядах з великою провідністю тривалість імпульсу набагато менша за скін-час (див.
~§~4.6), тож струм тече поверхнею плазми --- реалізується скінований пінч. При цьому тиск усередині
такого пінча має бути постійним, тому що \( j_z = 0 \), \( B = 0 \), а умова рівноваги
\eqref{eq:5.29} зводиться до рівності тиску плазми і магнітного тиску на межі.

У протилежному випадку, коли густина струму постійна по перерізу плазмового шнура \( j_z = I_0 / \pi
a^2 \), рівняння (5.18) визначає збільшення магнітного поля всередині плазми з радіусом
\begin{equation}\label{eq:5.30}
    B(r) = \frac{2I_0}{ca^2} r.
\end{equation}



З \eqref{eq:5.27} і \eqref{eq:5.30} знаходимо розподіл тиску по перерізу шнура:
\begin{equation}\label{eq:5.31}
    p(r) = \frac{I_0^2}{\pi c^2 a^2} \left( 1 - \frac{r^2}{a^2} \right),
\end{equation}
яке виявляється параболічним.


\end{document}